\section{Related work}
\label{sec:related_work}
\paragraph{Legend harmonization and multi-dataset labels.}
Land-cover harmonization addresses incompatible nomenclatures and mapping scales \citep{yang2017harmonization,baudoux2021translation}. MSeg additionally reconciles pixel annotations across datasets \citep{lambert2020mseg}; UniSeg formulates multi-dataset label conflicts and learns a shared representation \citep{kim2022uniseg}. The HTSS preprint integrates conflicting label spaces and different annotation types through a universal taxonomy \citep{meletis2023htss}. These works preclude treating taxonomy unification or heterogeneous-source learning as new here. Our operational audit records whether a mapped exact label actually occurs in the source masks used by the specified training protocol.

\paragraph{Open-set adaptation and recovery/retention reporting.}
Domain-adaptation methods address changes in representations and output distributions \citep{ganin2016dann,hoffman2018cycada,vu2019advent,zou2018classbalanced}. BUS and UniDA-SS explicitly handle differences between source and target class spaces and evaluate common/private performance \citep{choe2024bus,choe2025unida}. The SATS preprint also identifies known-class negative transfer and unknown-class underfitting \citep{ren2026sats}. Thus, balancing recovery and retention is established framing. The present estimands concern a named added leaf inside a fixed seven-leaf output space, with positive auxiliary-source supervision, rather than claiming a new open-set adaptation method or metric. We measure the tested mixtures' observed gain and cost without claiming superiority over these methods.

\paragraph{Transferable land-cover semantics and non-exact guidance.}
TACOSS uses textual supervision to align visual and semantic representations for transferable land-cover labels \citep{zermatten2025tacoss}. LandSegmenter combines heterogeneous land-use/land-cover supervision with flexible semantic representations \citep{liu2026landsegmenter}. These are close remote-sensing precedents for moving beyond a single rigid legend. Coarse or ambiguous guidance is also a learning problem in its own right \citep{li2024exactguidance,cour2011partial}. Our raw-mask support count is specific to the declared exact/coarse/ignored crosswalk; it does not establish that exact positive labels are universally necessary for learning or recognition. The source-mixture intervention changes imagery and labels jointly.


\paragraph{Recent semantic refinement and open-vocabulary segmentation.}
KG-OVRSeg enriches class representations with knowledge-graph relations to reduce naming and hierarchical ambiguity in remote-sensing segmentation \citep{huang2026kgovrseg}. Hierarchical Category Refinement studies the refinement of coarse land-cover categories using visual and semantic hierarchy information \citep{xiong2026hcr}. ROSS develops remote-sensing-specific representations and aggregation for text-defined open-vocabulary segmentation \citep{zhang2026ross}. These studies make the distinction between a software label assignment, semantic knowledge and botanical class identity especially important. Our support indicator concerns admitted source-mask labels; it neither measures all knowledge encoded by a pretrained model nor establishes that exact target-class annotations are necessary for recognition.

\paragraph{Coarse and historical product supervision.}
CTS-WS addresses temporal label changes and spatial-resolution discrepancies when existing products supervise finer-resolution mapping \citep{ctsws2025}. SAMPLER combines prototype learning and a modified Segment Anything model with label refinement for this coarse-to-fine setting \citep{sampler2025}. Superpixel-aware credible dual-expert learning instead uses superpixel training units and complementary expert estimates to handle noise in historical-product labels \citep{dualexpert2025}. These approaches address label reliability and resolution mismatch. Our audit asks a complementary question: which exact output assignments are actually admitted by the retained source-training protocol? An ignored or unmapped class and a noisy example of an admitted class are distinct supervision records. We do not benchmark these algorithms or import their reported gains into our source-mixture contrasts.

\paragraph{Domain generalization with foundation representations.}
Recent remote-sensing segmentation work adapts vision foundation models and expands source appearance diversity \citep{luo2025vfmdg}. MetaFD models domain variation continuously and disentangles domain-related features \citep{liang2025metafd}; concept-guided domain generalization learns semantic concept tokens while fine-tuning a pretrained vision foundation model for urban-scene segmentation in unseen domains \citep{liao2025concept}. Their focus on representations and appearance variation complements, rather than substitutes for, a source-label inventory. In our comparisons, auxiliary imagery and labels change together; these generalization mechanisms remain plausible contributors and cannot be isolated by a recovery score alone. The present CLIP and archived U-Net instruments therefore support bounded within-family analyses, not superiority over contemporary foundation-model methods.

\paragraph{Class-space adaptation and target-information requirements.}
Dual-classifier consistency and feature mixup have been studied for universal adaptation in remote-sensing classification \citep{li2025boostingunida}, while GLC++ addresses source-free universal adaptation through clustering and contrastive affinity learning \citep{qu2025glcpp}. Both provide recent category-shift precedents, but their classification/adaptation settings differ from our dense, fixed-output, positively supervised auxiliary-source experiment. One-shot remote-sensing segmentation adaptation also uses a target-image regime distinct from our source-mixture comparison \citep{tan2025oneshot}; ``one-shot'' here must not be read as ``open-set.'' We retain segmentation-specific common/private precedents above and make no new-metric or universal-adaptation claim.

\paragraph{Experimental instruments and scope.}
CLIP provides text-defined category scores \citep{radford2021clip}, and dense vision-language methods develop this interface for segmentation \citep{xu2022groupvit,zhou2022denseclip,li2022lseg}. The primary experiment freezes CLIP encoders and optimizes a low-rank calibration head. The archived ResNet-34 U-Net records provide a convolutional model-family check \citep{ronneberger2015unet}; its initialization and optimization regime differ from CLIP. LoveDA, FLAIR and OpenEarthMap define the empirical source/target setting \citep{wang2021loveda,garioud2023flair,xia2023openearthmap}. The paper reports within-family source contrasts, not a segmentation leaderboard or architecture-general performance comparison.
